% --- Archivo: 04_cuasi_newton.tex ---

\subsection{Métodos cuasi-Newton}

Una clase importante de métodos de optimización irrestricta es dada por el esquema iterativo
\begin{equation}
    x^{k+1} = x^k + \alpha_k d^k, \quad d^k = -Q_k f'(x^k), \quad k = 0, 1, \dots,
    \label{eq:cuasi_newton_esquema}
\end{equation}
donde, para todo $k$, $Q_k \in \R^{n \times n}$ es una matriz simétrica definida positiva, y $\alpha_k > 0$ es la longitud de paso calculada utilizando una de las reglas de búsqueda lineal. Notamos que el método es un método de descenso. Con efecto, si $x^k$ no es un punto estacionario del problema, como $Q_k$ es definida positiva, se tiene que
\begin{equation}
    \interno{f'(x^k)}{d^k} = -\interno{Q_k f'(x^k)}{f'(x^k)} < 0,
\end{equation}
y por lo tanto $d^k$ es una dirección de descenso de $f$ en $x^k$. En el caso de $Q_k = E^n$, el esquema se reduce a uno de los métodos del gradiente, y en el caso de $Q_k = (f''(x^k))^{-1}$ con $\alpha_k = 1$, obtenemos el método de Newton puro. Existen, sin embargo, otras posibilidades que son más atractivas computacionalmente.

\begin{algorithm}[Método de descenso con métrica variable]
    Elegir $x^0 \in \Rn$ y tomar $k := 0$. Elegir una de las reglas de búsqueda lineal y los parámetros asociados.
    \begin{enumerate}
        \item Elegir una matriz $n \times n$ simétrica definida positiva $Q_k$ y tomar $d^k = -Q_k f'(x^k)$. Calcular la longitud de paso $\alpha_k$ en la dirección $d^k$, por la regla de búsqueda lineal elegida (se $f'(x^k) = 0$, formalmente definimos $\alpha_k = 1$).
        \item Tomar $x^{k+1} = x^k + \alpha_k d^k$.
        \item Tomar $k := k + 1$ y retornar al Paso 1.
    \end{enumerate}
\end{algorithm}

Los métodos cuasi-Newton son métodos de esta forma, donde las matrices $Q_k$ son elegidas para aproximar, en un cierto sentido, la matriz $(f''(\bar{x}))^{-1}$ para alguna solución $\bar{x}$. Más precisamente, la secuencia $\{Q_k\}$ tiene que satisfacer la condición de Dennis-Moré, dada en el teorema siguiente.

\begin{theorem}[Teorema de Dennis-Moré]\label{thm:dennis_more}
    Supongamos que se satisfacen las hipótesis del Corolario 3.2.2. Además de eso, supongamos que la secuencia $\{x^k\}$, generada por el Algoritmo utilizando la regla de búsqueda lineal de Armijo con $\hat{\alpha} = 1$ e $\sigma \in (0, 1/2)$, converge a $\bar{x}$.
    
    Entonces, las siguientes dos afirmaciones son equivalentes:
    \begin{enumerate}[(a)]
        \item La tasa de convergencia de $\{x^k\}$ a $\bar{x}$ es superlineal, $\alpha_k = 1$ para todo $k$ suficientemente grande, y 
        \begin{equation}
            \|Q_k f'(x^k)\| \le M \|f'(x^k)\| \quad \forall k = 0, 1, \dots,
            \label{eq:dennis_more_a}
        \end{equation}
        donde $M > 0$ es una constante que no depende de $k$.
        
        \item Se tiene que
        \begin{equation}
            (Q_k - (f''(\bar{x}))^{-1})f'(x^k) = o(\|f'(x^k)\|).
            \label{eq:dennis_more_b}
        \end{equation}
    \end{enumerate}
\end{theorem}

\begin{proof}
    Primero, supongamos que la afirmación (a) se satisface. Llevando en cuenta que $\alpha_k = 1$ para todo $k$ suficientemente grande, tenemos que
    \begin{align}
        ((f''(\bar{x}))^{-1} - Q_k)f'(x^k) &= (f''(\bar{x}))^{-1}f'(x^k) + x^{k+1} - x^k \nonumber \\
        &= (f''(\bar{x}))^{-1}(f'(x^k) - f'(\bar{x}) - f''(\bar{x})(x^k - \bar{x})) + x^{k+1} - \bar{x} \nonumber \\
        &= o(\|x^k - \bar{x}\|),
        \label{eq:dem_dennis_1}
    \end{align}
    donde la última igualdad se sigue de la Fórmula de Taylor y de la convergencia superlineal de $\{x^k\}$ a $\bar{x}$. De \eqref{eq:dennis_more_a}, obtenemos que
    \begin{align*}
        M\|f'(x^k)\| &\ge \|Q_k f'(x^k)\| \\
        &\ge \|x^k - \bar{x}\| - \|x^{k+1} - \bar{x}\| = \|x^k - \bar{x}\| + o(\|x^k - \bar{x}\|).
    \end{align*}
    Por lo tanto,
    \[
        \frac{o(\|x^k - \bar{x}\|)}{\|f'(x^k)\|} \le \frac{o(\|x^k - \bar{x}\|)}{\|x^k - \bar{x}\|/M + o(\|x^k - \bar{x}\|)} \to 0 \quad (k \to \infty).
    \]
    Combinando esta última relación con \eqref{eq:dem_dennis_1}, obtenemos \eqref{eq:dennis_more_b}.

    Ahora supongamos que la afirmación (b) se satisface. La condición \eqref{eq:dennis_more_a} se sigue trivialmente (sino, obtenemos una contradicción obvia con \eqref{eq:dennis_more_b}). 
    Probaremos que $\alpha_k = 1$ para todo $k$ suficientemente grande. Por el Teorema del Valor Medio, existe $\tilde{t}_k \in [0, 1]$ tal que
    \[
        f(x^k - Q_k f'(x^k)) = f(x^k) - \interno{f'(x^k)}{Q_k f'(x^k)} + \frac{1}{2}\interno{f''(\tilde{x}^k)Q_k f'(x^k)}{Q_k f'(x^k)},
    \]
    donde $\tilde{x}^k = x^k - \tilde{t}_k Q_k f'(x^k)$. Para verificar que $\alpha_k = 1$ satisface la regla de Armijo, basta mostrar que
    \begin{equation}
        \interno{f'(x^k)}{Q_k f'(x^k)} - \frac{1}{2}\interno{f''(\tilde{x}^k)Q_k f'(x^k)}{Q_k f'(x^k)} \ge \sigma \interno{f'(x^k)}{Q_k f'(x^k)}.
        \label{eq:dem_dennis_2}
    \end{equation}
    Observamos que $\{\tilde{x}^k\} \to \bar{x}$ ($k \to \infty$), porque $\{x^k\} \to \bar{x}$ y $\{Q_k f'(x^k)\} \to 0$ (la última propiedad se sigue de \eqref{eq:dennis_more_a}). Por tanto, por \eqref{eq:dennis_more_b}, la desigualdad \eqref{eq:dem_dennis_2} puede ser escrita como
    \[
        (1 - \sigma)\interno{f'(x^k)}{(f''(\bar{x}))^{-1} f'(x^k)} \ge \frac{1}{2}\interno{f'(x^k)}{(f''(\bar{x}))^{-1} f'(x^k)} + o(\|f'(x^k)\|^2),
    \]
    o bien
    \[
        (1/2 - \sigma)\interno{f'(x^k)}{(f''(\bar{x}))^{-1} f'(x^k)} \ge o(\|f'(x^k)\|^2).
    \]
    Esta última desigualdad se satisface para todo $k$ suficientemente grande, porque $\sigma \in (0, 1/2)$, y el hecho de que $f''(\bar{x})$ es definida positiva implica que $(f''(\bar{x}))^{-1}$ es definida positiva, y por lo tanto, existe $\gamma > 0$ tal que
    \[
        (1/2 - \sigma)\interno{f'(x^k)}{(f''(\bar{x}))^{-1} f'(x^k)} \ge \gamma (1/2 - \sigma) \|f'(x^k)\|^2.
    \]
    Finalmente, probaremos que la tasa de convergencia de $\{x^k\}$ es superlineal. Para $k$ suficientemente grande, $\alpha_k = 1$, obtenemos
    \begin{align*}
        \|x^{k+1} - \bar{x}\| &= \|x^k - \bar{x} - Q_k f'(x^k)\| \\
        &= \|x^k - \bar{x} - (f''(\bar{x}))^{-1}f'(x^k) + ((f''(\bar{x}))^{-1} - Q_k)f'(x^k)\| \\
        &\le \|(f''(\bar{x}))^{-1}(f''(\bar{x})(x^k - \bar{x}) - f'(x^k))\| + o(\|f'(x^k)\|) \\
        &= o(\|x^k - \bar{x}\|),
    \end{align*}
    lo que concluye la demostración.
\end{proof}

Una de las conclusiones importantes del Teorema de Dennis-Moré es la siguiente: en la optimización de funciones dos veces diferenciables, para ser eficiente, un método debe hacer uso de aproximación de segunda orden de la función objetivo, de alguna manera. En particular, métodos que utilizan apenas aproximación de primer orden no pueden poseer convergencia superlineal en general. Cabe observar que para el método de Newton, la condición de Dennis-Moré se satisface automáticamente. 

No entanto, la información de segunda orden también puede ser construida implícitamente, sin cálculo directo de la derivada segunda de $f$. Esta es la idea principal de los métodos cuasi-Newton: la sustitución del cálculo de $f''(x^k)$ y de la resolución del sistema de ecuaciones lineales correspondiente, por una aproximación directa del paso calculada en cada iteración. Esto puede ser hecho de varias maneras diferentes. 

Para todo $k$, definimos
\begin{equation}
    r^k = x^{k+1} - x^k, \quad s^k = f'(x^{k+1}) - f'(x^k).
\end{equation}
Notamos que esta información ya está disponible en el momento de la escogencia de $Q_{k+1}$. Como queremos que esta matriz aproxime, en un cierto sentido, la inversa de la Hessiana de $f$, es natural imponer la condición
\begin{equation}
    Q_{k+1} s^k = r^k,
    \label{eq:ecuacion_secante}
\end{equation}
que es llamada de ecuación cuasi-Newton. Implícitamente admitiendo que la matriz $f''(\bar{x})$ es natural pedir la condición $s^k = Q_{k+1}^{-1} r^k$, que resulta en \eqref{eq:ecuacion_secante}.

Resumiendo, dada una matriz simétrica definida positiva $Q_k$ y vectores $r^k$ e $s^k$, la propuesta es escoger una matriz simétrica definida positiva $Q_{k+1}$ que satisfaga \eqref{eq:ecuacion_secante}. Históricamente, el primer método cuasi-Newton es el método de Davidon-Fletcher-Powell (DFP), que toma $Q_0$ como una matriz simétrica definida positiva arbitraria (digamos, $Q_0 = E^n$), y para todo $k$ define
\begin{equation}
    Q_{k+1} = Q_k + \frac{r^k (r^k)^\top}{\interno{r^k}{s^k}} - \frac{(Q_k s^k)(Q_k s^k)^\top}{\interno{Q_k s^k}{s^k}}.
    \label{eq:dfp_update}
\end{equation}
Notamos que todas las matrices así generadas serán simétricas. Además, la ecuación cuasi-Newton también será satisfecha.

\begin{proposition}\label{prop:dfp_pos_def}
    Sea $Q_k \in \R^{n \times n}$ una matriz simétrica definida positiva. Sea $f \colon \Rn \to \R$ una función diferenciable en los puntos $x^k, x^{k+1} \in \Rn$, que satisfacen la relación del paso con algún valor de longitud de paso $\alpha_k > 0$.
    Entonces las fórmulas definen sin ambigüedad una matriz $Q_{k+1}$ definida positiva si, y solamente si, la desigualdad
    \begin{equation}
        \interno{f'(x^{k+1})}{d^k} > \interno{f'(x^k)}{d^k}
        \label{eq:condicion_curvatura}
    \end{equation}
    es satisfecha.
\end{proposition}

\begin{proof}
    Supongamos que $Q_{k+1}$ sea definida positiva. Usando la ecuación cuasi-Newton, obtenemos
    \begin{align*}
        \interno{d^k}{f'(x^{k+1}) - f'(x^k)} &= \frac{1}{\alpha_k} \interno{r^k}{s^k} \\
        &= \frac{1}{\alpha_k} \interno{Q_{k+1} s^k}{s^k} > 0,
    \end{align*}
    donde la desigualdad vale suponiendo que $s^k \neq 0$ (caso contrario, $Q_{k+1}$ no estaría bien definida). Obtenemos entonces \eqref{eq:condicion_curvatura}.
    
    Supongamos ahora que \eqref{eq:condicion_curvatura} se satisfaga. Utilizando también el esquema obtenemos
    \[
        \interno{r^k}{s^k} = \alpha_k (\interno{d^k}{f'(x^{k+1}) - f'(x^k)}) > 0.
    \]
    En particular, $s^k \neq 0$ y, por el hecho de que $Q_k$ es definida positiva, $\interno{Q_k s^k}{s^k} > 0$. Esto muestra que $Q_{k+1}$ está bien definida. Para todo $\xi \in \Rn$, obtenemos
    \begin{align*}
        \interno{Q_{k+1}\xi}{\xi} &= \interno{Q_k \xi}{\xi} + \frac{\interno{r^k}{\xi}^2}{\interno{r^k}{s^k}} - \frac{\interno{Q_k s^k}{\xi}^2}{\interno{Q_k s^k}{s^k}} \\
        &= \frac{\|Q_k^{1/2} \xi\|^2 \|Q_k^{1/2} s^k\|^2 - \interno{Q_k^{1/2} \xi}{Q_k^{1/2} s^k}^2}{\|Q_k^{1/2} s^k\|^2} + \frac{\interno{r^k}{\xi}^2}{\interno{r^k}{s^k}} \ge 0,
    \end{align*}
    donde la desigualdad sigue del hecho de que ambos términos son no negativos por la desigualdad de Cauchy-Schwarz. Además, la igualdad a $0$ solamente puede valer cuando ambos términos se anulan, i.e., cuando $\interno{r^k}{\xi} = 0$ y $\|Q_k^{1/2} \xi\| \|Q_k^{1/2} s^k\| = |\interno{Q_k^{1/2} \xi}{Q_k^{1/2} s^k}|$. La última igualdad significa que $Q_k^{1/2} \xi = t (Q_k^{1/2} s^k)$ para algún $t$. Como $Q_k^{1/2}$ es no-singular, esto es posible cuando $t = 0$, i.e., $\xi = 0$. Por lo tanto, $Q_{k+1}$ es definida positiva.
\end{proof}

Como consecuencias del resultado arriba son las siguientes. Si el longitud de paso es calculado por la regla de minimización unidimensional, tenemos que $\interno{f'(x^{k+1})}{d^k} = 0$. Por tanto, valiendo automáticamente la condición, la relación de curvatura \eqref{eq:condicion_curvatura} siempre es satisfecha. Las reglas de Armijo y de Goldstein, no obstante, no poseen esta propiedad. Esta es la razón por la cual la regla de Wolfe es aquella utilizada con mayor frecuencia en implementaciones de métodos cuasi-Newton.

Entre todos los métodos cuasi-Newton, el método de Broyden-Fletcher-Goldfarb-Shanno (BFGS) es considerado actualmente el más eficiente. En este método, para todo $k$, definimos
\begin{equation}
    Q_{k+1} = Q_k + \frac{(r^k - Q_k s^k)(r^k)^\top + r^k(r^k - Q_k s^k)^\top}{\interno{r^k}{s^k}} - \frac{\interno{r^k - Q_k s^k}{s^k} r^k (r^k)^\top}{\interno{r^k}{s^k}^2}.
    \label{eq:bfgs_update}
\end{equation}
Es fácil ver que, como en el método DFP, las matrices generadas por \eqref{eq:bfgs_update} son simétricas, satisfacen a la ecuación cuasi-Newton, y la diferencia $Q_{k+1} - Q_k$ tiene puesto 2, como máximo.

Cuando $f$ es una función cuadrática dada y $Q$ es definida positiva, si el longitud de paso $\alpha_k$ es calculado por la regla de minimización unidimensional, los métodos DFP y BFGS poseen convergencia finita al único punto estacionario de $f$ (que en este caso es la única solución global del problema). En particular, el número de iteraciones $j$ necesario para encontrar la solución es menor o igual a la dimensión del espacio ($j \le n$) y se tiene que $Q_j = (f''(x^j))^{-1} = Q^{-1}$.